% Imporditud: tools/import_eksperimendid.py
% Allikas: Eesti füüsikaolümpiaadi eksperimendi ülesannete
%   kogu (Rael Kalda, 2023), ülesanne Ü131, lahendus L131.
\setAuthor{EFO žürii}
\setRound{lõppvoor}
\setYear{2005}
\setNumber{G E1}
\setDifficulty{4}
\setTopic{gaasid}
\setVahendid{Õhupall, millimeetripaber, marker, kaaluviht (veega täidetud plastpudel)}
\setPunktid{10}
\setAllikas{Eksperimendi kogu Ü131, lahendus lk 102}
\setLahendusseis{kasitsi}

\prob{Õhu tihedus}
Leidke õhupalli sees oleva õhu tihedus, kui on teada õhu temperatuur ja õhurõhk väljaspool palli. Hinnake mõõteviga. Võib lähtuda eeldusest, et palli väikeste deformatsioonide korral muutub palli kuju, mitte aga oluliselt tema ruumala. Õhu molaarmassiks lugege M = 29,0 kg/kmol. Õhurõhu ja temperatuuri väärtused tööruumis on toodud tahvlil.

\hint

\solu
\textit{Lahendus on kogumikus lk 102. Siia ei ole seda veel ümber kirjutatud, sest originaalis on tabel või joonis, mis PDF-i tekstikihist loetavalt välja ei tule.}
\probend
